(1) 自旋 \(1/2\) 粒子起始处于 \(S_z\) 的 \(+\hbar/2\) 本征态 \(|+z\rangle\)。求此态测量 \(S_x\) 的正、负结果概率。
(2) 若第一次 \(S_x\) 测量结果为正，随后测量 \(S_z\)，求两种结果概率。
(3) 若不筛选第一次 \(S_x\) 的测量结果，而将两路结果都保留，随后测量 \(S_z\)，求第二次 \(+z\) 的总概率。
(4) 推广到夹角为 \(\theta\) 的两测量轴 \(\hat n,\hat m\)，求从 \(+\hat n\) 得到 \(+\hat m\) 的条件概率并用Bloch矢量说明。
